% hw2-predicate-logic.tex

\newpage
\section{作业 2：谓词逻辑}

\begin{problem}[一阶谓词逻辑: 替换]
  请证明, 对于任意公式 $\alpha$ 与其中的任意变元 $x$,
  $x$ 都可以在 $\alpha$ 中无冲突地替换自己。

  \noindent (提示: 虽然该结论看上去很显然, 但请考虑如何以规范的方式证明它。
  你不能说 ``显然成立'' 或 ``不难看出'' 等等。)
\end{problem}

\begin{solution}
  证明命题:
  \[
    P(\alpha): \text{对任意变元 } x,\; x \text{ 都可以在 } \alpha \text{ 中无冲突地替换 } x.
  \]

  对公式 $\alpha$ 作结构归纳。

  \textbf{基础情况.}
  若 $\alpha$ 是原子公式, 则按 ``无冲突替换'' 的定义,
  原子公式情形无附加条件, 故对任意变元 $x$,
  $x$ 都可以在 $\alpha$ 中无冲突地替换 $x$。

  \textbf{归纳步骤.}

  \begin{enumerate}[(1)]
    \item 若 $\alpha = \lnot \psi$,
      且归纳假设表明 $P(\psi)$ 成立,
      则对任意变元 $x$,
      $x$ 可以在 $\psi$ 中无冲突地替换 $x$。
      由定义,
      $x$ 也可以在 $\lnot \psi$ 中无冲突地替换 $x$。

    \item 若 $\alpha = \psi \ast \gamma$,
      其中 $\ast \in \{\land, \lor, \to, \leftrightarrow\}$,
      且归纳假设表明 $P(\psi)$ 与 $P(\gamma)$ 都成立,
      则对任意变元 $x$,
      $x$ 在 $\psi$ 与 $\gamma$ 中都可以无冲突地替换自己。
      由定义,
      $x$ 也可以在 $\psi \ast \gamma$ 中无冲突地替换 $x$。

    \item 若 $\alpha = \forall y\; \psi$ 或 $\alpha = \exists y\; \psi$,
      固定任意变元 $x$。

      \begin{enumerate}[(a)]
        \item 若 $x = y$, 则 $x$ 在 $\forall/\exists x\; \psi$ 中不自由出现,
          故按定义中的条件 (a),
          $x$ 可以在 $\forall/\exists x\; \psi$ 中无冲突地替换 $x$。

        \item 若 $x \neq y$, 则变元 $y$ 不在代入项 $x$ 中出现,
          且由归纳假设, $x$ 可以在 $\psi$ 中无冲突地替换 $x$。
          因而按定义中的条件 (b),
          $x$ 可以在 $\forall/\exists y\; \psi$ 中无冲突地替换 $x$。
      \end{enumerate}
  \end{enumerate}
  综上, $P(\alpha)$ 对任意公式 $\alpha$ 都成立。
\end{solution}
%%%%%%%%%%%%%%%

\begin{problem}[一阶谓词逻辑: 前束范式]
  对于任意的命题逻辑公式,
  我们在上次作业中介绍了如何将其化归成合取范式或析取范式。
  类似地, 对于任意的一阶谓词逻辑公式,
  我们也可以将其化归成前束范式 (prenex normal form),
  即``所有量词都出现在公式最前面''的形式。

  准确地说, 称具有以下形式的公式为前束范式:
  \[
    Q_1 x_1\; Q_2 x_2\; \cdots\; Q_n x_n\; \alpha
  \]
  其中, 每个 $Q_i$ 都是 $\forall$ 或 $\exists$,
  并且 $\alpha$ 中无量词。

  可以利用下列规则进行量词操作, 将任意的一阶谓词逻辑公式化归成前束范式:
  \begin{enumerate}[label = \textbf{R\arabic*}]
    \item $\lnot \forall x\; \alpha \equiv \exists x\; \lnot \alpha$;
    \item $\lnot \exists x\; \alpha \equiv \forall x\; \lnot \alpha$;
    \item $\alpha \to \forall x\; \beta \equiv \forall x\; (\alpha \to \beta)$,
      其中 $x$ 不在 $\alpha$ 中自由出现;
    \item $\alpha \to \exists x\; \beta \equiv \exists x\; (\alpha \to \beta)$,
      其中 $x$ 不在 $\alpha$ 中自由出现;
    \item $(\forall x\; \alpha) \to \beta \equiv \exists x\; (\alpha \to \beta)$,
      其中 $x$ 不在 $\beta$ 中自由出现;
    \item $(\exists x\; \alpha) \to \beta \equiv \forall x\; (\alpha \to \beta)$,
      其中 $x$ 不在 $\beta$ 中自由出现.
  \end{enumerate}

  请利用以上规则将以下公式化归成前束范式,
  并给出每一步使用的规则:
  \[
    \forall x\; \exists y\; P(x, y) \to \exists x\; Q(x).
  \]

  (提示: 答案可能不唯一)
\end{problem}

\begin{solution}
  先对右侧量词中的束缚变元换名:
  \[
    \forall x\; \exists y\; P(x, y) \to \exists x\; Q(x)
    \equiv
    \forall x\; \exists y\; P(x, y) \to \exists z\; Q(z).
  \]

  下面开始化归:
  \begin{align*}
    \forall x\; \exists y\; P(x, y) \to \exists z\; Q(z)
      &\equiv \exists x\; \big(\exists y\; P(x, y) \to \exists z\; Q(z)\big)
        && \text{(R5; $x$ 不在 $\exists z\; Q(z)$ 中自由出现)} \\
      &\equiv \exists x\; \forall y\; \big(P(x, y) \to \exists z\; Q(z)\big)
        && \text{(R6; $y$ 不在 $\exists z\; Q(z)$ 中自由出现)} \\
      &\equiv \exists x\; \forall y\; \exists z\; \big(P(x, y) \to Q(z)\big)
        && \text{(R4; $z$ 不在 $P(x, y)$ 中自由出现)}.
  \end{align*}

  因而可取其前束范式为
  \[
    \boxed{\exists x\; \forall y\; \exists z\; \big(P(x, y) \to Q(z)\big)}.
  \]
\end{solution}